% SPDX-License-Identifier: GPL-3.0-or-later
% Copyright (C) 2026 Alessandro Greco
\chapter{Difese contro DoS e gestione delle risorse}\label{cap:15-dos}
\aggiornatoal{23 settembre 2026}{torspec commit \texttt{8e1f245}; tor 0.4.9.13; dos-spec, hspow-spec}

\begin{obiettivi}
Al termine di questo capitolo saprai:
\begin{itemize}[nosep]
  \item spiegare perché un attacco di negazione del servizio contro Tor è anche un attacco all'anonimato;
  \item descrivere come tor protegge memoria, CPU e banda dei relay;
  \item conoscere i limiti per indirizzo su connessioni e circuiti e il ruolo dei parametri di consenso;
  \item spiegare come gli onion service si difendono dalle inondazioni di richieste di introduzione;
  \item descrivere la prova di lavoro di tipo v1 (Equi-X e Blake2b) e il suo controllo dello sforzo;
  \item provocare e osservare le difese di un relay e di un onion service nella rete di laboratorio (Laboratorio~9).
\end{itemize}
\end{obiettivi}

Una rete aperta, anonima e gestita da volontari è un bersaglio naturale per chi vuole renderla inutilizzabile. In questo capitolo vediamo come Tor si difende dagli attacchi di \emph{negazione del servizio} (\emph{denial of service}, DoS): tentativi di esaurire le risorse di un relay o di un onion service per impedirne l'uso. Gli episodi concreti di questi anni, come l'ondata di attacchi agli onion service del 2022--2023, sono raccontati nel \cref{cap:27-casi-rete}; qui ci concentriamo sui meccanismi.

% ---------------------------------------------------------------------------
\section{Perché il DoS è un problema di anonimato}\label{sec:dos-anonimato}

In una rete qualsiasi, un attacco DoS è \enquote{solo} un problema di disponibilità. In Tor è anche un'arma contro l'anonimato~\spec{dos-spec}{Overview}:
\begin{itemize}
  \item un avversario che riesce a spegnere a comando un relay può costringere i client che lo usano come guard a sceglierne un altro, sperando che sia suo. Jansen, Tschorsch, Johnson e Scheuermann hanno mostrato come esaurire la memoria di un relay con poche risorse e usare questa capacità per deanonimizzare gli onion service: il \emph{Sniper Attack}~\cite{jansen2014sniper};
  \item rendere lenti o irraggiungibili i relay onesti spinge il traffico verso quelli dell'avversario, che ne vede una quota maggiore;
  \item un attacco che fa fallire circuiti in modo mirato può essere usato come segnale, per esempio per capire quali circuiti passano da un certo relay.
\end{itemize}
Per questo le difese non si limitano a \enquote{tenere in piedi} i relay: cercano di farlo senza creare, a loro volta, comportamenti osservabili e sfruttabili da un avversario.

% ---------------------------------------------------------------------------
\section{Le risorse fisiche}\label{sec:dos-risorse}

\subsection{Memoria}

La memoria è la risorsa più delicata, e quella protetta dai meccanismi più maturi~\spec{dos-spec}{Memory exhaustion}. Su Linux, un programma che chiede troppa memoria di solito non riceve un errore: il sistema gliela promette e poi, quando la usa davvero, termina il processo. Aspettare l'errore di allocazione non funziona. Tor si impone quindi un limite, \opt{MaxMemInQueues}, impostato dall'operatore o calcolato come frazione della memoria del sistema, e tiene il conto di ciò che occupa: celle in coda sui circuiti, buffer delle connessioni, stati di compressione, descrittori degli onion service in cache, risoluzioni DNS in cache.

Quando il totale supera il limite, tor libera memoria in due passi: prima svuota le cache che occupano più del 20\,\% del totale, partendo dalle voci più vecchie, finché non scendono al 10\,\%; poi, se è ancora oltre il 90\,\% del limite, chiude circuiti e connessioni ordinandoli in base all'età del dato più vecchio che tengono in coda, cominciando da quelli con i dati fermi da più tempo. L'idea è colpire chi sta accumulando dati senza leggerli, cioè proprio il comportamento tipico di un attacco.

\subsection{CPU}

La CPU è il punto debole dell'implementazione in C, che elabora la rete su un unico thread principale. Le operazioni più costose, gli handshake di circuito (\cref{cap:09-circuiti}), sono delegate a un gruppo di thread di lavoro e messe in coda; ma se ne arrivano troppe, la coda si allunga e i circuiti legittimi scadono. Tor non misura direttamente l'uso della CPU: la protegge indirettamente, limitando le operazioni che la consumano, come la creazione di circuiti e le introduzioni agli onion service.

\subsection{Banda}

Gli operatori possono limitare la banda con opzioni come \opt{RelayBandwidthRate}, ma questi limiti servono a non superare il proprio contratto, non a difendersi. Contro un uso anomalo della banda Tor si affida al controllo di flusso e della congestione (\cref{cap:11-stream}) e a una distribuzione equa delle celle tra i circuiti.

% ---------------------------------------------------------------------------
\section{Limiti per indirizzo}\label{sec:dos-limiti}

La difesa più efficace contro gli attacchi grossolani è semplice: un guard vede l'indirizzo IP di chi gli si collega, e può limitare quanto ciascun indirizzo può chiedere. È l'unico punto della rete in cui questo è possibile, perché più avanti nel circuito l'origine del traffico è, per costruzione, sconosciuta. Tor applica tre famiglie di limiti ai client (non ai relay autenticati), configurabili con opzioni \texttt{DoS\ldots}:
\begin{description}
  \item[Connessioni] (\texttt{DoSConnection\ldots}) un massimo di connessioni TLS contemporanee e un ritmo massimo di nuove connessioni per indirizzo; oltre la soglia le connessioni vengono chiuse.
  \item[Circuiti] (\texttt{DoSCircuitCreation\ldots}) un \emph{secchio di gettoni} (\emph{token bucket}) per indirizzo: ogni nuovo circuito consuma un gettone, i gettoni si ricaricano a un ritmo fisso fino a un massimo (il \emph{burst}). Un indirizzo che ha esaurito il secchio \emph{e} tiene aperte almeno un certo numero di connessioni viene marcato, e per un periodo i suoi nuovi circuiti vengono rifiutati.
  \item[Rendezvous a un salto] (\texttt{DoSRefuseSingleHopClientRendezvous}) il rifiuto delle richieste di rendezvous che arrivano su circuiti di un solo salto, come quelli che usavano i gateway Tor2web per raggiungere gli onion service senza anonimato per il client: un modo economico di generare carico per la rete.
\end{description}

\subsection{Il codice e il consenso}

Nel codice di tor queste difese sono \textbf{disattivate} per impostazione predefinita (\file{DOS\_CC\_ENABLED\_DEFAULT 0} in \file{src/core/or/dos.h}); a regolarle, per tutta la rete, sono i parametri di consenso. Il consenso pubblico del 23 settembre 2026 contiene:
\begin{center}\small
\begin{tabular}{lrr}
\toprule
Parametro & Consenso & Predefinito nel codice \\
\midrule
\texttt{DoSCircuitCreationEnabled} & 1 & 0 \\
\texttt{DoSCircuitCreationMinConnections} & 2 & 3 \\
\texttt{DoSCircuitCreationRate} & 2 & 3 \\
\texttt{DoSCircuitCreationBurst} & 60 & 90 \\
\texttt{DoSConnectionEnabled} & 1 & 0 \\
\texttt{DoSConnectionMaxConcurrentCount} & 50 & 100 \\
\texttt{DoSRefuseSingleHopClientRendezvous} & 1 & --- \\
\bottomrule
\end{tabular}
\end{center}
I valori della rete sono più severi di quelli del codice: sono il frutto di anni di attacchi e di regolazioni successive. Poterli cambiare dal consenso, in poche ore e senza aggiornare il software di migliaia di relay, è uno degli strumenti più preziosi del Tor Project durante un attacco.

\begin{attenzione}[Il limite dei limiti per indirizzo]
Un avversario con molti indirizzi IP, come una botnet, sfugge ai limiti per indirizzo; e un limite troppo severo colpisce gli utenti legittimi che condividono lo stesso indirizzo, per esempio dietro il NAT di un operatore mobile o di un'università. I valori sono sempre un compromesso.
\end{attenzione}

% ---------------------------------------------------------------------------
\section{Onion service sotto attacco}\label{sec:dos-onion}

Gli onion service sono particolarmente vulnerabili. Ogni richiesta di introduzione costa pochissimo al client: un \texttt{INTRODUCE1} (\cref{sec:os-rendezvous}). Al servizio invece costa molto: decifrare la richiesta, eseguire un handshake e costruire un circuito di più salti verso il punto di rendezvous indicato. È un'\textbf{amplificazione}: con poche risorse un avversario costringe il servizio, e tutti i relay dei circuiti che costruisce, a un lavoro molto maggiore. E il servizio non può bloccare gli indirizzi degli attaccanti, perché non li vede.

\subsection{Limiti al punto di introduzione}

Una prima difesa è far limitare le richieste ai punti di introduzione. Con l'opzione \texttt{HiddenService\-EnableIntroDoS\-Defense} e le varianti \texttt{\ldots RatePerSec} e \texttt{\ldots BurstPerSec}, il servizio comunica a ciascun punto di introduzione, con un'estensione di \file{ESTABLISH\_INTRO} (\file{DOS\_PARAMS}), quante \texttt{INTRODUCE2} al secondo è disposto a ricevere. Il limite protegge il servizio, ma non distingue un attaccante da un utente legittimo: sotto attacco, entrambi vengono respinti.

\subsection{La prova di lavoro}

La difesa più recente, proposta con la \proposta{327} e disponibile da tor 0.4.8.4 (agosto 2023), è la \textbf{prova di lavoro} (\emph{proof of work}, PoW)~\spec{hspow-spec}{Motivation}. L'idea è permettere ai client di dimostrare di aver speso una certa quantità di calcolo e di dare precedenza alle richieste che ne dimostrano di più. \textbf{Non è una soglia minima obbligatoria}: lo sforzo suggerito dal servizio è un'indicazione per il client, non un requisito d'accesso; le introduzioni senza prova sono ammesse in coda con priorità zero quando la difesa è abilitata. Una prova \emph{presente ma non valida}, invece, è rifiutata. Il costo per il client dipende dallo sforzo scelto e dall'hardware; per un attaccante che vuole inviare richieste in grande quantità il costo si moltiplica.

Il meccanismo ha quattro parti~\spec{hspow-spec}{Common protocol}:
\begin{enumerate}
  \item \textbf{I parametri nel descrittore.} Il servizio pubblica, nel secondo strato del descrittore (\cref{sec:os-descrittore}), una riga \texttt{pow-params} con lo schema (\texttt{v1}), un \emph{seme} casuale di 32 byte, lo \emph{sforzo suggerito} e la scadenza del seme, rinnovato ogni 105--120 minuti.
  \item \textbf{Il client risolve il problema.} Con lo sforzo $E$ scelto, il client cerca una soluzione del problema Equi-X per la sfida
  \[ \texttt{"Tor hs intro v1\textbackslash 0"} \,\|\, A' \,\|\, \mathit{seme} \,\|\, \mathit{nonce} \,\|\, E, \]
  dove $A'$ è la chiave accecata del servizio. Una soluzione Equi-X non basta: il client calcola un hash Blake2b di 32 bit, $R$, della sfida e della soluzione, e la soluzione è accettabile solo se $R \cdot E \le 2^{32} - 1$. Altrimenti cambia nonce e riprova.
  \item \textbf{La prova nell'introduzione.} Il client inserisce nonce, sforzo, i primi 4 byte del seme e la soluzione in un'estensione della parte cifrata di \texttt{INTRODUCE1}, che solo il servizio legge.
  \item \textbf{La coda a priorità.} Il servizio verifica la prova, che è molto più veloce da verificare che da trovare, e mette la richiesta in una coda ordinata per sforzo. Serve per prime le richieste con lo sforzo più alto; quelle senza prova hanno sforzo zero. Se la coda si riempie, scarta di colpo la metà con lo sforzo più basso.
\end{enumerate}

\begin{figure}[htbp]
\centering
% SPDX-License-Identifier: GPL-3.0-or-later
% Copyright (C) 2026 Alessandro Greco
%% Testo alternativo per l'accessibilita' (usato solo nell'EPUB).
\altfig{Protezione PoW facoltativa per servizi onion. Se il servizio la abilita, pubblica pow-params con seme e sforzo suggerito E; il client sceglie l'effort effettivo, cerca nonce e soluzione Equi-X che soddisfano anche R per E non maggiore di 2^32−1, e inserisce la prova nell'estensione cifrata di INTRODUCE1. Il punto di introduzione inoltra INTRODUCE2 senza leggere la prova, con eventuale limite al ritmo. Il servizio verifica seme, nonce, Equi-X e test di sforzo, scartando prove non valide. Le richieste valide entrano nella coda ordinata per sforzo verificato decrescente; senza prova l'effort è zero. Quando la coda supera il limite, l'implementazione può eliminare la metà con sforzo inferiore. I valori 900, 400, 120, 40 e 0 sono esempi, non soglie fisse.}
% La prova di lavoro degli onion service: dal descrittore alla coda a priorità
\begin{tikzpicture}[font=\sffamily\scriptsize,
  box/.style={draw,rounded corners=2pt,align=center,minimum height=1cm}]
  \node[box,fill=torgrigio!8,text width=2.2cm] (d) at (0,0) {descrittore:\\\texttt{pow-params}\\seme, sforzo suggerito};
  \node[box,fill=torviola!10,text width=3.2cm] (c) at (3.8,0) {client: sceglie $E$, cerca nonce\\e soluzione Equi-X con\\$R \cdot E \le 2^{32}-1$};
  \node[box,fill=torgrigio!8,text width=2.2cm] (i) at (8.4,0) {punto di\\introduzione\\(limite al ritmo)};
  \node[box,fill=torverde!12,text width=2.2cm] (s) at (8.4,-2.2) {servizio:\\verifica la prova\\(veloce)};
  \draw[->] (d) -- (c);
  \draw[->] (c) -- node[below,pos=0.45] {\texttt{INTRODUCE1}} (i);
  \draw[->] (i) -- node[right] {\texttt{INTRODUCE2}} (s);
  \foreach \k/\e/\c in {0/900/torverde!45,1/400/torverde!35,2/120/torverde!25,3/40/torverde!15,4/0/torgrigio!10}
    \node[draw,minimum width=1.1cm,minimum height=0.55cm,fill=\c] (q\k) at (1.2+1.1*\k,-2.2) {$E{=}\e$};
  \draw[->] (s) -- (q4.east);
  \node[above=0.1cm of q2] {coda ordinata per sforzo (senza prova: $E=0$)};
  \node[below=0.1cm of q0,align=center] {servita per prima};
  \node[below=0.1cm of q4,align=center] {coda piena: si scarta\\la metà con sforzo minore};
\end{tikzpicture}

\caption{La prova di lavoro per gli onion service. Il client legge seme e sforzo suggerito dal descrittore, trova una soluzione che supera il test, la invia nella parte cifrata di \texttt{INTRODUCE1}; il servizio la verifica e ordina le richieste per sforzo.}
\label{fig:pow}
\end{figure}

\begin{approfondimento}[Perché Equi-X]
Equi-X, progettato da tevador per questo scopo, è un problema \emph{asimmetrico}: trovare una soluzione richiede memoria e molte operazioni, verificarla pochissime. Si basa su HashX, una famiglia di funzioni di hash che cambia implementazione a ogni seme e sfrutta l'esecuzione in parallelo delle CPU moderne: è difficile da accelerare con hardware dedicato, così che un avversario non guadagni molto rispetto a un utente con un normale computer. Esistono due implementazioni, quella in C inclusa in tor e una in Rust scritta per Arti, tenute compatibili con test automatici. Il codice C è distribuito con licenza LGPL: per includerlo, tor va compilato con l'opzione \texttt{--enable-gpl}, e un tor compilato senza di essa non offre la PoW: \cmd{tor --list-modules} riporta allora \texttt{pow: no} invece del \texttt{pow: yes} che vedremo nel \cref{cap:17-torrc}. Non è però il caso comune: sono compilati con l'opzione sia l'immagine di laboratorio sia il pacchetto \texttt{tor} di Debian, e in entrambi il modulo risulta presente.
\end{approfondimento}

\subsection{Uno sforzo lineare}

Perché il livello Blake2b? Il costo di una singola soluzione Equi-X è grosso modo fisso. Il test $R \cdot E \le 2^{32}-1$ è superato da una soluzione casuale con probabilità circa $1/E$: il client deve quindi provare in media $E$ soluzioni. Lo sforzo cresce \emph{linearmente} con $E$, e questo rende semplice regolarlo. Molti sistemi di prova di lavoro chiedono invece un certo numero di bit a zero, con un costo che raddoppia a ogni bit: una scala troppo grossolana per un controllo fine. Lo sforzo è inoltre dichiarato \emph{prima} di cercare la soluzione e fa parte della sfida: un client non può sperare in un colpo di fortuna e dichiarare poi uno sforzo più alto di quello speso.

\subsection{Il controllo dello sforzo}

Nessuno fissa lo sforzo una volta per tutte. Il servizio parte da uno sforzo suggerito di zero, cioè con la difesa dormiente, e ogni cinque minuti valuta lo stato della coda~\spec{hspow-spec}{Common protocol}. Aumenta il valore interno se ha scartato richieste con sforzo \emph{strettamente maggiore} di quello suggerito oppure se la coda si è congestionata e contiene ancora richieste con sforzo \emph{almeno pari} al suggerimento. Lo diminuisce soltanto quando nessuna delle due condizioni vale e la coda è abbastanza breve; altrimenti lo lascia invariato. Nell'aumento sceglie il massimo fra \emph{valore precedente più uno} e il rapporto fra sforzo totale accodato e rendezvous effettivamente avviati; nella diminuzione moltiplica per $2/3$. È una logica ispirata al controllo della congestione di TCP, ma aggiornata a intervalli fissi. Non ogni variazione del valore interno comporta la ripubblicazione del descrittore. I client, a loro volta, possono aumentare lo sforzo dopo la scadenza di un tentativo di collegamento. È un sistema a più anelli di retroazione, nel quale un client può scegliere di investire uno sforzo maggiore del valore suggerito per ottenere priorità.

% ---------------------------------------------------------------------------
\section{Laboratorio: provocare le difese}\label{sec:lab-dos}

\begin{laboratorio}[9]{Limiti anti-DoS e prova di lavoro nella rete di laboratorio}
\textbf{Obiettivo.} Attivare i limiti sulla creazione di circuiti di un relay e superarli; attivare la prova di lavoro sull'onion service e leggerne i parametri; studiare il livello Blake2b della prova di lavoro con \file{labs/it/codice/torlab/pow.py}.\\
\textbf{Prerequisiti.} Laboratori~1--8; la rete di laboratorio avviata, con l'immagine compilata con \texttt{--enable-gpl}.\\
\textbf{Tempo stimato.} 2 ore.
\end{laboratorio}

\begin{attenzione}[Solo sulla tua rete]
Questo laboratorio genera di proposito traffico anomalo. Va eseguito esclusivamente sulla rete privata del libro: fare lo stesso contro relay o servizi reali è un attacco, dannoso per altre persone e, in molti ordinamenti, illegale (\cref{cap:25-legale}).
\end{attenzione}

\paragraph{Passo 1: attivare i limiti su un relay.} Aggiungiamo al \file{torrc} del relay \texttt{test003r} limiti molto bassi, per vederli scattare subito, e facciamogli rileggere la configurazione con il segnale \texttt{HUP}:
\begin{lstlisting}[language=torrc,numbers=none]
DoSCircuitCreationEnabled 1
DoSCircuitCreationMinConnections 2
DoSCircuitCreationRate 1
DoSCircuitCreationBurst 10
DoSCircuitCreationDefenseType 2
DoSCircuitCreationDefenseTimePeriod 60 seconds
\end{lstlisting}
Il tipo di difesa 2 significa \enquote{rifiuta i nuovi circuiti}.

\paragraph{Passo 2: superarli.} Apriamo uno o due canali verso il relay e chiediamo 30 circuiti in rapida successione con \cella{CREATE\_FAST}. Per avere contatori separati in ogni prova, facciamo partire le connessioni da indirizzi di loopback diversi (127.0.0.2, 127.0.0.3), con il nuovo parametro \texttt{origine} di \file{apri\_canale}:
\begin{lstlisting}[numbers=none]
origine 127.0.0.2, 2 canali: 9 ok, 21 DESTROY(5); primo rifiuto al n. 10
origine 127.0.0.3, 1 canale:  30 ok
origine 127.0.0.1, 2 canali: 9 ok, 21 DESTROY(5); primo rifiuto al n. 10
\end{lstlisting}
Tre osservazioni:
\begin{itemize}
  \item il relay rifiuta con \cella{DESTROY} di motivo 5, \texttt{RESOURCELIMIT};
  \item con un \emph{burst} di 10 passano 9 circuiti, non 10: il codice (\file{cc\_has\_exhausted\_circuits} in \file{src/core/or/dos.c}) marca il client appena il secchio arriva a zero, quindi la decima richiesta è già respinta;
  \item con un solo canale non scatta nulla, perché la difesa richiede almeno \texttt{MinConnections} connessioni contemporanee. Il motivo è evitare di punire un normale client, che di solito tiene aperta una sola connessione al proprio guard.
\end{itemize}
Avevamo previsto che le connessioni da 127.0.0.1, l'indirizzo dei relay del laboratorio, fossero esenti dai limiti; l'esperimento ha smentito la previsione, e il codice di \file{dos.c} conferma che un'esenzione del genere non esiste per i circuiti. Alla fine dell'esperimento ripristiniamo il \file{torrc} originale.

\paragraph{Passo 3: la prova di lavoro dell'onion service.} Aggiungiamo, dopo le righe \texttt{HiddenService\ldots} del \file{torrc} di \texttt{test009h}:
\begin{lstlisting}[language=torrc,numbers=none]
HiddenServicePoWDefensesEnabled 1
\end{lstlisting}
Dopo il segnale \texttt{HUP}, il servizio pubblica un nuovo descrittore. Scaricandolo e decifrandolo con il Laboratorio~8 troviamo, nel secondo strato:
\begin{lstlisting}[numbers=none]
pow-params v1 tx2Mtmzw5GHgWjcDoivTvWzwQPNZONwMwSMddOc+ZVQ= 0 2026-09-23T18:20:53
\end{lstlisting}
Schema \texttt{v1}, il seme in base64, uno sforzo suggerito di \textbf{zero}, perché il servizio non ha ancora visto una coda, e la scadenza del seme, un'ora e cinquanta minuti dopo la pubblicazione, nella finestra di 105--120 minuti prevista.

\paragraph{Passo 4: lo sforzo è lineare.} Il modulo \file{pow.py} implementa \textbf{soltanto il controllo Blake2b} e usa byte casuali come sostituti statistici dei candidati Equi-X: consente di studiare la distribuzione del valore di controllo \emph{assumendo che l'hash sia uniforme}, non di produrre una prova valida per Tor. Non esegue \texttt{equix\_solve} né \texttt{equix\_verify}, non controlla la scadenza del seme e non gestisce la cache anti-replay. La prova reale richiede \emph{tutti} questi passaggi~\spec{hspow-spec}{Onion service proof-of-work: Scheme v1, Equi-X and Blake2b}:
\begin{lstlisting}[language=Python]
def valore_controllo(id_servizio, seme, nonce, sforzo, soluzione):
    sfida = P_STRING + id_servizio + seme + nonce + struct.pack(">I", sforzo)
    h = hashlib.blake2b(sfida + soluzione, digest_size=4).digest()
    return int.from_bytes(h, "big")

def sforzo_valido(r, sforzo):
    return r * sforzo <= UINT32_MAX
\end{lstlisting}
Misurando in media quante soluzioni servono per 400 prove a ciascuno sforzo:
\begin{lstlisting}[numbers=none]
sforzo    1: 1.0 tentativi
sforzo   10: 9.7
sforzo  100: 102.3
sforzo 1000: 1068.8
\end{lstlisting}
Il numero di tentativi cresce in proporzione allo sforzo, come previsto: il valore atteso è lo sforzo stesso, perché ogni tentativo riesce con probabilità $1/\mathit{sforzo}$ e il numero di tentativi segue una distribuzione geometrica. Di quella distribuzione lo scarto quadratico medio è grande quanto la media, quindi la media di 400 prove ha uno scarto tipico del $1/\sqrt{400} = 5\,\%$: fra due esecuzioni i valori possono differire del dieci per cento senza che nulla sia cambiato, e rieseguendo la misura abbiamo infatti ottenuto 969,9 dove qui si legge 1068,8. È un buon promemoria: una media va sempre letta insieme all'incertezza del campione. Si noti poi un dettaglio: \texttt{blake2b} con \file{digest\_size=4} non equivale ai primi 4 byte di un Blake2b di 64 byte, perché la lunghezza dell'uscita entra nello stato iniziale della funzione. La specifica lo sottolinea, ed è un errore facile da commettere.

% ---------------------------------------------------------------------------
\section*{Riepilogo}
\begin{itemize}
  \item In Tor gli attacchi DoS sono anche attacchi all'anonimato: spegnere relay a comando o farli fallire può spostare il traffico verso l'avversario.
  \item La memoria è protetta da un limite autoimposto, \opt{MaxMemInQueues}, oltre il quale tor libera cache e chiude i circuiti che accumulano dati.
  \item I guard limitano per indirizzo connessioni e circuiti dei client con secchi di gettoni; le difese, spente nel codice, sono attivate e regolate dal consenso.
  \item Gli onion service subiscono un'amplificazione: introdursi costa poco, rispondere molto.
  \item La prova di lavoro v1 combina Equi-X e un test Blake2b che rende lo sforzo lineare; il servizio ordina le richieste per sforzo e regola il valore suggerito in base alla coda. Il valore suggerito non è una soglia obbligatoria, e il laboratorio simula soltanto il test Blake2b.
\end{itemize}

% ---------------------------------------------------------------------------
\section*{Esercizi}

\begin{esercizio}[1]\label{es:15-bucket}
Con \texttt{DoSCircuitCreationRate} 2 e \texttt{Burst} 60, come nel consenso pubblico, quanti circuiti può creare un client in un minuto partendo da un secchio pieno? E in un'ora? Supponi che la difesa sia già attiva, cioè che l'indirizzo abbia raggiunto \texttt{DoSCircuitCreationMinConnections} connessioni contemporanee: sotto quella soglia non c'è alcun limite.
\end{esercizio}

\begin{esercizio}[1]\label{es:15-guard}
Perché i limiti per indirizzo si applicano solo ai client, e solo sul primo salto? Che cosa succederebbe se un middle cercasse di limitare i circuiti \enquote{per indirizzo}?
\end{esercizio}

\begin{esercizio}[2]\label{es:15-amplificazione}
Stima il rapporto tra il lavoro del client e quello della rete per una singola richiesta di introduzione: quanti circuiti, handshake e celle genera la risposta del servizio?
\end{esercizio}

\begin{esercizio}[2]\label{es:15-bit}
Confronta la prova di lavoro di Tor con uno schema che richiede $k$ bit iniziali a zero in un hash. Quali valori di sforzo può esprimere ciascuno schema? Perché la scala lineare è preferibile per un controllo a retroazione?
\end{esercizio}

\begin{esercizio}[2]\label{es:15-sforzo}
Un client dichiara uno sforzo $E = 50$ ma, per fortuna, trova subito una soluzione con $R$ molto piccolo, che sarebbe valida anche per $E = 5000$. Può inviarla dichiarando $E = 5000$? Perché?
\end{esercizio}

\begin{esercizio}[3]\label{es:15-oom}
Nella rete di laboratorio, imposta \opt{MaxMemInQueues} al valore minimo su un relay intermedio e costruisci un circuito attraverso di esso da cui scarichi un file grande \emph{senza leggere} i dati in arrivo (e senza inviare SENDME). Che cosa succede alle celle accumulate? Cerca nei log del relay i messaggi del gestore della memoria.
\end{esercizio}

\begin{esercizio}[3]\label{es:15-coda}
Scrivi un piccolo simulatore della coda a priorità di un onion service: arrivano richieste con sforzo casuale (molte a zero, poche alte), il servizio ne serve un certo numero al secondo e scarta la metà più bassa quando la coda supera una soglia. Misura il tempo d'attesa medio in funzione dello sforzo.
\end{esercizio}

% ---------------------------------------------------------------------------
\section*{Per approfondire}
\begin{itemize}
  \item \titolo{Tor denial-of-service mitigations} (dos-spec) e \titolo{Onion service proof-of-work} (hspow-spec)~\cite{torspec}.
  \item R.~Jansen, F.~Tschorsch, A.~Johnson, B.~Scheuermann, \titolo{The Sniper Attack: Anonymously Deanonymizing and Disabling the Tor Network}~\cite{jansen2014sniper}.
  \item Il repository di Equi-X di tevador e il suo diario di sviluppo, citati nella specifica.
\end{itemize}
